\scanpage{50}
% E: text|upper left, question|J'aimerais comprendre
J'aimerais comprendre la relation entre les \(\sigma_a^F\), et les \uline{signes} \(\varepsilon_F(a)\).

% E: text|upper left, transition|Mais il faut encore
Mais il faut encore introduire une autre structure.

% E: text|left upper-middle, c|Ensemble C_K
c) Ensemble \(C_K\unsure{\subset|\in}S_K\) des ``pts critiques''\inserted{ils sont d'ordre 2, \unsure{tels que}}, avec la condition que sur chaque \deleted{\unsure{arc}} grand côté, il y a \deleted{\unsure{au plus}}\inserted{exactement} un \deleted{tel} élément de \(C_K\).

\deleted{(quand il n'y en a pas, l'arête est dite exceptionnelle)} De plus, chaque pt critique \(s\) est muni d'un poids \(\nu(s)\geq\unsure{2|1}\),\inserted{qui est pair ou impair, suivant que \(s\) est noir, ou rouge,} (le poids d'un élément de \(S_K\) est par définition \(\nu(s)\) si critique, 1 si non critique \deleted{\unsure{de poids}} et dans \(S'_K\), 0 s'il est non critique et est dans \(S''_K\)).

\margin{définir les ``arcs critiques'' ou ``arcs de transition'' d'amplitude totale sur les côtés, \(\nu(s)\) ou \unsure{encore} \(2\nu(s)\) (il contient \unsure{\(\nu(s)\) subdivisions} \deleted{de \(S_K\)}\inserted{\unsure{\(S_2\)}} \unsure{et sommets rouges}).}
\margin{Les pts critiques sont des centres des \unsure{arcs critiques}, et déterminent les indices (\unsure{comme la moitié de l'amplitude des arcs critiques dans le cas}) --- \unsure{les voisins des sommets sur ces arcs sont non critiques}.}
\margin{\uline{NB} Les arcs de transition (fermés) sont mutuellement disjoints.}
\margin{Les sommets d'ordre 4 sont exactement \deleted{\unsure{ceux qui sont les}} les \unsure{milieux des arcs} \deleted{\unsure{antipodiques}}\inserted{critiques} \deleted{\unsure{\(\forall\) les arcs antipodiques}}.}

% E: text|left lower-middle, introduction|En termes des sommets critiques
En termes des sommets critiques, on a ceci :

% E: text|left lower, 1|Soit s un sommet critique
1) Soit \(s\) un sommet critique, incident à la face \(F\). Soit \(s'\) son antipodique dans \(\delta F\). \deleted{Alors si \(s'\) est aussi critique, \unsure{on doit le signe \(\varepsilon_a\)}}\inserted{Pour qu'on ait \(\varepsilon_F(a)=+1\), il f. et il s.} que les sommets\inserted{\(s'\)} antipodiques de \(s\)\inserted{\unsure{critiques}} soient également critiques (il a alors le $=$ indice).
\margin{relation entre \(C_K\) et les signes \(\varepsilon_F(a)\)}

% E: text|bottom left, 2 unfinished|Pour toute face F
2) Pour tte face \(F\), il y a au moins 3 pts critiques sur \(\delta F\) (et $=$ au moins 4, \deleted{\unsure{3q si tous les pts}} il y a au moins un pt critique\inserted{\(s\)} de signe \(+1\). \deleted{\unsure{Pour tout}} \unsure{cas, sur chacun}
\mspage{49}
